\documentclass[12pt]{article}
\usepackage[utf8]{inputenc}
\usepackage[margin=0.75in]{geometry}
\usepackage{amsmath}
\usepackage{amssymb}

\title{CSC110 Fall 2026 Assignment 2}
\author{Idzudin Saffuan}
\date{\today}

\begin{document}
\maketitle

% TIP: In LaTeX, the underscore (_) is a special character, so if you want to use it
% in normal text, you have to put a backslash in front of it. E.g., a2\_part2.py,
% not a2_part2.py.

\section*{Part 1: Proofs}

\begin{enumerate}
\item[1a.]
Let $x \in \mathbb{Z}^+ $

Take $y=x$ , hence $y \in \mathbb{Z}^+$

\textcircled{1} $xy \implies x^2 \geq 0$

\textcircled{2} $ y =x \implies x = x \implies y \geq x$



\item[1b.]
WTS: $\forall{y} \in D, \exists{x} \in D\ s.t \ xy \leq 0 \vee y < x $

Let $ y \in \mathbb{Z}^+$

\textcircled{1} Take $x = 2y$ , hence $ x \in \mathbb{Z}^+$

\textcircled{2} $ y < x \implies y < 2y$

\item[1c.]
Suppose $D = \mathbb{Z}^-$

Let $x \in \mathbb{Z}^-$

\textcircled{1} Take $y = 2x$

\textcircled{2} Hence, the statement $y \geq x \implies 2x \geq x$ is false.

\item[1d.]
Suppose $D = (0,1]$

Let $x \in D$

Take $y = 1$

\textcircled{1} $xy \implies x > 0 $ since $y=1$ and $x \in (0,1]$

\textcircled{2} Hence $\forall{x}, \ y \geq x$ since $0 < x \leq 1 = y$



\item[2.]
WTS: $a>1  \land  b>1 \implies (\exists{d}\in\mathbb{N},(\exists{k}\in\mathbb{Z},\ ab = kd))\land(d\neq1 \land d\neq ab)$

Suppose $a>1$ and $b>1$

\textcircled{1} Take $d = b$ where $(d\neq1 \land d\neq ab)$

\textcircled{2} Take $k = a$

\textcircled{3} Hence $ab = ab$


\item[3.]
Suppose $a \in \mathbb{Z}$

Then by Quotient-Remainder Theorem,

$a = 2q + r$ for some $q,r \in \mathbb{Z}$ and $0 \leq r \leq 1$

$a^2 + a \implies (2q+r)^2 + 2q +r$

$= 4q^2 +4qr + r^2 + 2q+r$

$= 4q^2 +4qr + r(4q+1+r)$

\underline{Case 1: $r=1$;}

$4q^2 + 2q + (4q+1+1)$

$=4q^2 + 6q +2$

$= 2(2q^2 +3q +1)$

\underline{Case 2: $r=0$;}

$4q + 2q$

$=2(2q^2+q)$

For both of these cases, it is shown that $2 \ | \  a^2 +a$.




\item[4a.]
Suppose $S= \{0,1,2,3,4\}, \ \forall a_1, a_2, a_3, b_1, b_2, b_3 \in S, \forall r \in \mathbb{Z}, \mathit{AssessDiff}(a_1, a_2, a_3, b_1, b_2, b_3, r)$

It is given that $r = r_1 +r_2 + r_3$

Hence $r = |a_1-b_1| + |a_2 - b_2| + |a_3 - b_3|$

Since $a_i, b_i \in S \ for \ i \in \{1,2,3\},$

$|a_i - b_i| \geq 0$, since absolute values are always non-negative

$max\{S\} = 4$, and $min\{s\} = 0$

Hence absolute difference between any two elements in S is;

$|4-0| = 4$

Thus $r_i \leq 4$

Combining these ranges gives:

$0 \leq r_i \leq 4$




\item[4b.]
Suppose $S= \{0,1,2,3,4\}, \ \forall a_1, a_2, a_3, b_1, b_2, b_3 \in S, \forall r \in \mathbb{Z}, \mathit{AssessDiff}(a_1, a_2, a_3, b_1, b_2, b_3, r)$

\textcircled{1} Assume that $r =0$

Hence $0 = |a_1-b_1| + |a_2 - b_2| + |a_3 - b_3|$

for $|a_n - b_n| = 0$, $a_n = b_n$ for $n$ in \{1,2,3\}

\textcircled{2} Now, Assume instead $a_n = b_n$ for $n$ in \{1,2,3\}

Then $r = |a_1-b_1| + |a_2 - b_2| + |a_3 - b_3|$

$= 0 + 0+ 0$

Hence $r=0$.




\item[4c.]
Suppose $S= \{0,1,2,3,4\}, \ \forall a_1, a_2, a_3, b_1, b_2, b_3 \in S, \forall r \in \mathbb{Z}, \mathit{AssessDiff}(a_1, a_2, a_3, b_1, b_2, b_3, r)$

Then $r_n = |a_n-b_n|$:

$=|(-1)(b_n-a_n)|$

$=|(-1)||b_n-a_n|$

$= 1|b_n-a_n|$

$|b_n-a_n|$





\item[4d.]
Let $S' = \{0,1,2,3,4,5\}$

WTS: \ $\exists a_1, a_2, a_3, b_1, b_2, b_3 \in S', \exists r \in \mathbb{Z}, \mathit{AssessDiff}(a_1, a_2, a_3, b_1, b_2, b_3, r)$ and $r>12$

Take $a_i = 0 \ and\  b_i = 5 \ for \ i \in \{1,2,3\}$

Hence r = 15
\end{enumerate}

\section*{Part 2: Loops and mutation debugging}

\begin{enumerate}
\item[1.]

test\_routine\_report PASSED

test\_monitor\_report FAILED

test\_urgent\_report FAILED



\item[2.]

Error 1: The function clean\_text returns texts that are only split but not all lowercase. This is due to strings being immutable.

Error 2: The function count\_severity\_words only checks if the given word is in the list
occurrences\_so\_far and adds it to the list and given an occurrences of +1. But in the case where the given word is already in the list, it does not add 1 to the word occurrences in the list.

\item[3.]
test\_routine\_report passed as the severity words in the document occur only once and all are in lowercase letters.
\end{enumerate}

\end{document}
